\begin{table*}[t]
\centering
\footnotesize
\setlength{\tabcolsep}{3pt}
\begin{tabular}{lrrrr}
\toprule
Population & $n$ & mean cov. & worst & $<0.85$ \\
\midrule
Buildings (BDG2), metered & 18 & \textbf{0.901} & 0.824 & 1 \\
System demand, metered & 6 & \textbf{0.829} & 0.762 & 5 \\
System demand, reconstructed & 6 & \textbf{0.824} & 0.565 & 2 \\
\midrule
\emph{India, one city (I-BLEND + Delhi SLDC), windows} &  &  &  &  \\
\quad buildings, native 15-min & 12 & \textbf{0.910} & 0.897 & 0 \\
\quad campus feed (one HT consumer) & 1 & \textbf{0.873} & 0.873 & 0 \\
\quad city (Delhi SLDC) & 1 & \textbf{0.762} & 0.762 & 1 \\
\bottomrule
\end{tabular}
\caption{Empirical coverage of the nominal 90\% interval after conformal calibration, \emph{as the benchmark calibrated it}: split conformal plus an adaptive layer at the step $\gamma=0.35$ in the units of each series. Read this way the guarantee very nearly holds on individual buildings and fails systematically on system-level demand. The third row is the replication: six Chinese provinces, a different country, a different year and a different data-generating process, reproduce the system-level failure to within 0.005 of the metered panel and contain the worst row in the study (Heilongjiang, 0.565).}
\label{tab:calibration}
\end{table*}
